

\begin{enumerate}
	\item 知乎： 关于超螺旋滑模算法的完整证明过程 \url{https://zhuanlan.zhihu.com/p/672355794}
	\item CSDN： 超螺旋滑模控制详细介绍（全网独家） \url{https://blog.csdn.net/qq_41811966/article/details/134869940}
\end{enumerate}

\subsection{稳定性条件概览}

为了尽可能地兼顾一般性和行文简便，这里将\eqref{eq sta s}写为：
\begin{equation}
	\begin{cases}
		\dot{s}=-k_1 \sig{s}^{1/2}+w+\varrho_1\\
		\dot{w}=-k_2 \sgn{s}+\varrho_2
	\end{cases}
\end{equation}
这里$\varrho_1,\varrho_2$分别为第一通道和第二通道的扰动。
事实上，这里的第一通道扰动可以合并到第二通道，
即记$w'=w+\varrho_1$这样$\dot{w}'=\dot{w}+\dot{\varrho_1}=-k_2 \sgn{s}+\varrho_2+\dot{\varrho_1}$.
于是，系统可以转化为$\varrho_1'=0$, $\varrho_2'=\varrho_2+\dot{\varrho}_1$.
因此，表\ref{table sta gain selection}给出了已有文献对增益选择的要求。

\begin{table}[!htb]
	\centering
	\begin{tabular}{c|c}
		\hline
		参数条件  & 出处 
		\\\hline  
		\makecell{
		当 $\varrho_2=0$ 时，$k_1>0$并且 $k_2>0$ ，\\
		当 $\varrho_2>0$ 时，通过算法\ref{algo sta moreno}确定
		}
		&
			\makecell{
			引自\cite[algorithm 1]{morenoStrictLyapunovFunctions2012}，\\
			\cite[第五章]{YangGaoJieHuaMoKongZhiLiLunJiQiZaiQianQuDongXiTongZhongDeYingYongYanJiu2016}
		}
		\\\hline
		$k_1 >0$,
		$k_2>\|\dot{\omega}\|_\infty$& 
		引自多体文献\cite{fengFinitetimeDistributedConvex2020}，
		有误
		\\\hline
		$k_1 >\sqrt{6} k_2$,
		$k_2>\|\dot{\omega}\|_\infty$&
		改写自多体文献\cite{chenDistributedFinitetimeDifferentiator2024}，
		不够保守
		\\\hline
	\end{tabular}
	\caption{STA收敛条件对比，在第一通道无扰动$\varrho_1=0$的参数对比}
	\label{table sta gain selection}
\end{table}

\begin{algo}\label{algo sta moreno}
	\begin{enumerate}
		\item  选择正常数$(\beta,\gamma)$ 使得 $0<\beta<1,\gamma>1$
		\item 找到满足下列不等式的正常数 $\chi,\alpha$
		\begin{equation}
			\chi-\frac{2}{\gamma}\alpha
			>\alpha^2-\beta (1+\chi)\alpha+\frac14 (1+\chi)^2
		\end{equation}
		\item 选取增益参数为
		\begin{equation}
			k_1=\chi \sqrt{\frac{2\gamma}{(1-\beta)\alpha}}\sqrt{L},\quad 
			k_2=\frac{\beta+1}{1-\beta}L
		\end{equation}
	\end{enumerate}
\end{algo}